\section{Curling-curls}
\label{app:curling-curls}

\apref{app:curling-curls} fixes names and compact count forms. Cageon confinement is developed in Section~\ref{sec:cageon-confinement}.

\subsection{Naming conventions}
\label{app:cc-naming-conventions}

\subsubsection{Hinge triad}
\label{app:cc-hinge-triad}
\[
\mathsf H_3=(x^-,\mu,x^+).
\]
Here \(x^-\) is the reflected branch, \(\mu\) is the hinge / dwell / pivot, and \(x^+\) is the successor branch. The local ternary support is
\[
\{-1,0,+1\}.
\]

\subsubsection{Curl}
\label{app:cc-curl}
Curl is returned closure through a hinge; in the main text it is typed as Stokes-aligned curl motif / closure witness.

\subsubsection{Nested curl}
\label{app:cc-nested-curl}
A nested curl is a curl carried inside another curl.

\subsubsection{Pin-curl}
\label{app:cc-pin-curl}
A pin-curl is a returned curl whose residue seats, delays, or fixes a hinge.

\subsubsection{Curling curls}
\label{app:cc-curling-curls}
Curling curls are recursive returns that become active path structure inside the fractal tesseract topology.

\subsubsection{Monon}
\label{app:cc-monon}
\[
\mathrm{monon}(x)=\operatorname{cycle}(x\succ x=x).
\]

\subsubsection{Dion}
\label{app:cc-dion}
\[
D_\alpha=(+1,-1),
\qquad
D_\Omega=(-1,+1).
\]

\subsubsection{Duon, duad, wavelet}
\label{app:cc-duon-duad-wavelet}
\[
\mathrm{duon}=D_i\mu D_j,
\qquad
D_i,D_j\in\{D_\alpha,D_\Omega\},
\qquad
\mu\in\{-1,0,+1\}.
\]
The resolved family has
\[
2\times3\times2=12
\]
variants. The open-seat duon is \(E_2^0\), and the seat-closed duon is the duad \(E_2^{\pm}\).
\[
\mathrm{wavelet}=W_\omega[\mathrm{duon}],
\qquad
W_\omega^\Delta[\mathrm{duon}],
\qquad
\Delta=\Delta^{\chi\mathrm{lag}}.
\]

\subsubsection{Odd core primitives}
\label{app:cc-odd-core-primitives}
The odd core primitives are monon and trion. They are complete cycle/knot primitives and cageable cores.

\subsubsection{Tetrion, tetron, tetratrion}
\label{app:cc-tetrion-tetron-tetratrion}
\[
\mathrm{tetrion}=\text{bare four-cycle / four-pivot return closure}.
\]
\[
\mathrm{tetron}=E_4=\text{default inner cage with four seats}.
\]
\[
\mathrm{tetratrion}=4\times\mathrm{trion}=3:4.
\]

\subsubsection{Core families}
\label{app:cc-core-families}
\[
1:2=\mathrm{Dimonon}=2\times\mathrm{monon},
\qquad
\operatorname{stem}(1:2)=\mathrm{Dimon}.
\]
\[
3:1=\mathrm{Trio}=1\times\mathrm{trion},
\qquad
3:2=\mathrm{Ditrion}=2\times\mathrm{trion}.
\]
\[
3:3=\mathrm{Tritrion}=3\times\mathrm{trion},
\qquad
\operatorname{stem}(3:3)=\mathrm{Tritrio}.
\]
\[
3:4=\mathrm{Tetratrion}=4\times\mathrm{trion},
\qquad
\operatorname{stem}(3:4)=\mathrm{Tetratrio}.
\]

\subsubsection{Outer cage names}
\label{app:cc-even-cages}
\[
\mathsf C_6=\mathrm{hexon},
\qquad
\mathsf C_8=\mathrm{octon},
\qquad
\mathsf C_{10}=\mathrm{decon}.
\]

\subsubsection{Standalone caged names}
\label{app:cc-standalone-caged-names}
Standalone caged names combine the inner family stem with the outer cage name.
\[
1:2:6=\mathrm{Dimonhexon},
\qquad
3:1:6=\mathrm{Triohexon},
\qquad
3:2:6=\mathrm{Ditriohexon}.
\]
\[
3:3:6=\mathrm{Tritriohexon},
\qquad
3:4:6=\mathrm{Tetratriohexon}.
\]
\[
1:2:8=\mathrm{Dimonocton},
\qquad
1:2:10=\mathrm{Dimondecon}.
\]

\subsection{Count compression}
\label{app:cc-count-compression}

\subsubsection{Count basis}
\label{app:cc-count-basis}
\[
\mathcal B_{\mathrm{curl}}=(3:3,
3:4,
1:2)=(\mathrm{Tritrio},
\mathrm{Tetratrio},
\mathrm{Dimon}).
\]
The compact count form is
\[
a:b:c::L,
\]
where \(a=\#(3:3)\), \(b=\#(3:4)\), \(c=\#(1:2)\), and \(L\) is outer cage order.

\subsubsection{Basic shared hexon}
\label{app:cc-basic-shared-hexon}
\[
1:1:0::6
=
\mathrm{Tritrio}::\mathrm{Tetratrio}::\mathrm{hexon}.
\]
Expanded:
\[
(3:3)^1::(3:4)^1::(1:2)^0::6.
\]

\subsubsection{Filled shared cages}
\label{app:cc-filled-shared-cages}
\[
3:3:0::6
=
\mathrm{Tritrio}^{3}::\mathrm{Tetratrio}^{3}::\mathrm{hexon}.
\]
\[
4:4:0::8
=
\mathrm{Tritrio}^{4}::\mathrm{Tetratrio}^{4}::\mathrm{octon}.
\]
\[
5:5:0::10
=
\mathrm{Tritrio}^{5}::\mathrm{Tetratrio}^{5}::\mathrm{decon}.
\]
General filled shared cage:
\[
n:n:0::2n
=
\mathrm{Tritrio}^{n}::\mathrm{Tetratrio}^{n}::\operatorname{out}(2n),
\qquad
n=3,4,5,\ldots .
\]

\subsubsection{Tetrad facilitation}
\label{app:cc-tetrad-facilitation}
Tetratrio may function as the full tetrad seat for Tritrio:
\[
1:1:0::6
=
\mathrm{Tritrio}::\mathrm{Tetratrio}::\mathrm{hexon}.
\]

\subsection{Active naming set}
\label{app:cc-active-naming-set}

\subsubsection{Primitives}
\label{app:cc-active-primitives}
\[
\mathsf H_3=\text{hinge triad}.
\]
\[
\mathrm{monon},
\quad
\mathrm{dion},
\quad
\mathrm{duon},
\quad
\mathrm{duad}.
\]
\[
\mathrm{tetrion}=4\text{-cycle},
\qquad
\mathrm{tetron}=E_4=\text{default inner }4\text{-cage}.
\]

\subsubsection{Core families}
\label{app:cc-active-core-families}
\[
1:2=\mathrm{Dimonon},
\qquad
3:1=\mathrm{Trio},
\qquad
3:2=\mathrm{Ditrion},
\]
\[
3:3=\mathrm{Tritrion},
\qquad
3:4=\mathrm{Tetratrion}.
\]

\subsubsection{Caged names}
\label{app:cc-active-caged-names}
\[
1:2:6=\mathrm{Dimonhexon},
\qquad
3:3:6=\mathrm{Tritriohexon},
\qquad
3:4:6=\mathrm{Tetratriohexon}.
\]
\[
1:2:8=\mathrm{Dimonocton},
\qquad
1:2:10=\mathrm{Dimondecon}.
\]

\subsubsection{Shared cage names}
\label{app:cc-active-shared-cage-names}
\[
1:1:0::6
=
\mathrm{Tritrio}::\mathrm{Tetratrio}::\mathrm{hexon}.
\]
\[
3:3:0::6
=
\mathrm{Tritrio}^{3}::\mathrm{Tetratrio}^{3}::\mathrm{hexon}.
\]
